% Generated draft: requires mathematical and source review.
\documentclass[11pt]{article}
\usepackage[margin=25mm]{geometry}
\usepackage{fontspec}
\setmainfont{Latin Modern Roman}
\usepackage{amsmath,amssymb,mathtools,unicode-math}
\setmathfont{Latin Modern Math}
\usepackage{enumitem,needspace,xurl,hyperref}
\setlength{\parindent}{0pt}
\setlength{\parskip}{5pt}
\setlength{\emergencystretch}{3em}
\newcommand{\N}{\mathbb{N}}
\newcommand{\Z}{\mathbb{Z}}
\newcommand{\Q}{\mathbb{Q}}
\newcommand{\R}{\mathbb{R}}
\newcommand{\C}{\mathbb{C}}
\newcommand{\F}{\mathbb{F}}
\newcommand{\A}{\mathbb{A}}
\newcommand{\PP}{\mathbb{P}}
\newcommand{\E}{\mathbb{E}}
\newcommand{\eps}{\varepsilon}
\newcommand{\dd}{\,\mathrm{d}}
\begin{document}
\section*{UnsolvedMath ID: 10800018}
\textbf{Source problem code:} catalogue entry
\subsection{Definitions and mathematical statement}
\emph{Problem (Victor Vassiliev, Problem 7).} Give a similar universal estimate for the radius of convergence of the multidimensional Newton method of [Shub and Smale1993]. The preceding estimate to which ``similar'' refers is not included in the statement.

In the conventional analytic square-system setting, let \(f:\Omega\subseteq\mathbb K^m\to\mathbb K^m\), \(\mathbb K=\R\) or \(\C\), have a zero \(x^*\) with invertible derivative. Newton iteration is
\[
 x_{j+1}=x_j-Df(x_j)^{-1}f(x_j)
\]
whenever the inverse exists. Relative to a specified norm and domain, a guaranteed convergence radius \(r\) means that every \(x_0\) with \(\|x_0-x^*\|<r\) has all iterates defined in the domain and converging to \(x^*\). The request is for a universal lower estimate for such a radius in terms of the intended data. The original data, allowed function class, field, norm, domain, convergence rate, and dependence on dimension or degree are not supplied; this conventional model does not resolve those missing choices.

The research summary attributes a solution to Shub--Smale \(\alpha\)-theory (1993), describes later refinements and implementations including alphaCertified and interval-arithmetic certification, and proposes the numerical bound \(1/(2\gamma(f,x^*))\) or an \(\alpha\)-criterion threshold, analogous to an earlier \(d/(2n-1)\) estimate. In a usual analytic convention,
\[
 \gamma(f,x^*)=\sup_{k\geq2}
 \left\|Df(x^*)^{-1}\frac{D^k f(x^*)}{k!}\right\|^{1/(k-1)}.
\]
The quoted numerical constant and its required hypotheses are not established by the supplied record and are not asserted as a universal theorem here. The summary does not define the \(d,n\) of its comparison, state a precise \(\alpha\)-criterion, or reconcile these data with the source's intended meaning of ``similar.
\subsection{Short English statement}
Seek a universal guaranteed convergence neighborhood for multidimensional Newton iteration, analogous to an omitted preceding estimate. The registry labels the problem solved by \(\alpha\)-theory, but its proposed numerical \(\gamma\)-bound and the intended comparison require verification.
\subsection{Sources}
\begin{itemize}
\item Local UnsolvedMath registry, permanent ID 10800018, source code catalogue entry; proposer Victor Vassiliev, Problem 7.
\item Vassiliev, \emph{A Few Problems on Monodromy and Discriminants}, Problem 7, with the supplied citation [Shub and Smale1993]. Full bibliographic details of that citation are not provided. Supplied article URL: \url{https://amj.math.stonybrook.edu/html-articles/Files-2015-2024/15-11/}; recorded access 2026-07-29; extraction html-structure. The article contains eighteen explicitly labeled theorem-style problem environments.
\item The supplied record has status solved and SOLVED-IN-LITERATURE, crediting Shub--Smale \(\alpha\)-theory and subsequent certification work. Its generic background instead describes the source as presenting an open problem and retains NEEDS\_REVIEW for a current primary-literature audit. The specific numerical bound and claimed resolution are supplied assessments, not independently verified facts. Redistribution terms for the publicly accessible author/publisher copy also require review. The registry assigns L2 whereas background records default L3.
\end{itemize}
\end{document}
